\section{Representaciones de Grupos Finitos}

\begin{ejerciciostar}
    Demostrar que $T_q\in SO(V)$ para todo cuaternio $q$ de norma 1.
\end{ejerciciostar}

\begin{ejerciciostar}\label{ej:3_2}
    Calcular explícitamente una representación real no trivial de grado 2 del grupo de permutaciones $S_3$.
\end{ejerciciostar}

\begin{ejercicio}
    Comprobar que la multiplicación definida sobre $KG$ es asociativa. Su elemento neutro es $1e$, donde $e$ es el elemento neutro de $G$.
\end{ejercicio}

\begin{ejemplo}\label{ejm:ejs3_4}
    para cada $q\in \bb{H}$ tal que $\|q\|^2 = qq^\ast = 1$, consideramos la aplicación $T_q:V\to V$ dada por $T_q(v) = qvq^\ast$, donde $V$ es el subespacio vectorial real de $\bb{H}$ con base $\{i,j,k\}$. Que $T_q(v)\in V$ para todo $v\in V$ es consecuencia de que, según vimos en la demostración del Teorema 1.8.2, $V = \{a\in \bb{H}: a^2 \leq 0\}$.

    Por otra parte, resulta fácil demostrar que $T_q$ es $\mathbb{R}-$lineal y que preserva la norma de $V$, por lo que $T_q\in O(V)$. De hecho, $T_q\in SO(V)$ (es decir, es un giro), lo que no es tan inmediato, pero no es demasiado difícil de ver. En cualquir caso, $T_q\in GL(V)$. Además $T_{pq} = T_p\circ T_q$, para todo $p,q\in \bb{H}$ de norma 1.

    Consideremos ahora $\cc{Q}=\{\pm 1, \pm i, \pm j, \pm k\}$ que es un subgrupo multiplicativo de $\bb{H}^\times$ con 8 elemento, llamado \textit{grupo cuaterniónico}. Por todo lo discutido, la aplicación $T:\cc{Q}\to GL(V)$, definida por $T(q)= T_q$ para $q\in \cc{Q}$, proporciona una representación real de $\cc{Q}$ de grado (o dimensión) 3.
\end{ejemplo}

\begin{ejercicio}
    Calcular todos los subespacios invariantes para la representación de $\cc{Q}$ del Ejemplo~\ref{ejm:ejs3_4}.
\end{ejercicio}

\begin{ejercicio}
    Calcular todos los subespacios invariantes para la representación de $S_3$ del Ejercicio~\ref{ej:3_2}.
\end{ejercicio}

\begin{coro}\label{coro:ejs6_7}
    Sean $(V_1,\rho_1), \ldots, (V_t,\rho_t)$ las representaciones irreducibles complejas de $G$, con multiplicidades $(n_1,\ldots,n_t)$ en la representación regular. Entonces $\dim_\mathbb{C} V_i = n_i$ para $i = 1,\ldots,t$ y $|G| = n_1^2 + \ldots + n_t^2$.
\end{coro}

\begin{ejercicio}
    Deducir del Corolario~\ref{coro:ejs6_7} cuántas representaciones irreducibles complejas no equivalentes tiene $S_3$, y cuáles son sus dimensiones. Calcularlas explícitamente.
\end{ejercicio}

\begin{ejercicio}
    Deducir del Corolario~\ref{coro:ejs6_7} cuántas representaciones irreducibles complejas no equivalentes tiene $\cc{Q}$, y cuáles son sus dimensiones. Calcularlas explícitamente.
\end{ejercicio}

\begin{ejercicio}
    Sea $H$ un subgrupo normal de $G$ y $\pi:G\to G/H$ la proyección canónica. Demostrar que $(V,\rho)$ es una representación irreducible de $G/H$ si, y sólo si, $(V,\rho\pi)$ es una representación irreducible de $G$.
\end{ejercicio}

\begin{ejercicio}
    Demostrar que, para todo $a\in \mathbb{C}G$, se tiene que $\widetilde{\chi_\rho}(a) = tr\tilde{\rho}(a)$.
\end{ejercicio}

\begin{ejercicio}
    Demostrar que todo carácter de un grupo finito $G$ es constante sobre cada clase de conjugación de $G$.
\end{ejercicio}

\begin{ejercicio}
    Calcular la tabla de caracteres del grupo cíclico $C_4$.
\end{ejercicio}

\begin{ejerciciostar}
    Calcular razonadamente la tabla de caracteres del grupo $\cc{Q}$ definido en el Ejemplo~\ref{ejm:ejs3_4}.
\end{ejerciciostar}

\begin{ejercicio}
    Los caracteres irreducibles de $S_4$ proporcionan, por restricción, algunos caracteres de $A_4$. Describir estos caracteres. ¿Son irreducibles?
\end{ejercicio}

\begin{ejerciciostar}
    Calcular razonadamente la tabla de caracteres del grupo diédrico $D_4$.
\end{ejerciciostar}

\begin{ejerciciostarstar}
    Calcular razonadamente la tabla de caracteres del grupo diédrico $D_n$, para $n\geq 2$.
\end{ejerciciostarstar}

\begin{ejercicio}\label{ej:3_16}
    Sea $G$ un grupo abeliano finito, y sea $\hat{G}$ el conjunto de los caracteres complejos irreducibles de $G$. Demostrar que el producto inducido por el de números complejos dota a $\hat{G}$ de estructura de grupo.
\end{ejercicio}

\begin{ejerciciostarstar}
    Sea $G$ un grupo abeliano finito, y $\hat{G}$ el grupo definido en el Ejercicio~\ref{ej:3_16}. Demostrar que existe un isomorfismo de grupos $G\cong \hat{G}$.
\end{ejerciciostarstar}

\begin{ejercicio}
    Sea $G$ un grupo finito y $g\in G$. Demostrar que $g$ es conjugado con $g^{-1}$ si, y sólo si, $\chi(g)\in \mathbb{R}$ para todo carácter complejo irreducible $\chi$ de $G$.
\end{ejercicio}

\begin{ejercicio}
    Demostrar que $\varphi^G \in \cc{C}(G)$ para cada $\varphi \in \cc{C}(H)$, y que
    \begin{equation*}
        \varphi^G(1) = [G:H]\varphi(1)
    \end{equation*}
\end{ejercicio}

\begin{ejerciciostar}
    Sea $G$ un grupo abeliano finito, y $H$ un subgrupo de $G$. Demostrar que la aplicación $(-)_H:\cc{C}(G)\to \cc{C}(H)$ es sobreyectiva. Identificar su núcleo.
\end{ejerciciostar}

\begin{ejerciciostar}
    Sea $\bb{F}_5$ el cuerpo de 5 elementos y
    \begin{equation*}
        G = \left\{\left(\begin{array}{cc}
            a & b \\
            0 & a^{-1} 
        \end{array}\right) : a,b\in \bb{F}_5, a\neq 0\right\}
    \end{equation*}
    Comprobar que $G$ es un subgrupo de $GL_2(\bb{F}_5)$ y calcular razonadamente la tabla de caracter de $G$.
\end{ejerciciostar}

\begin{ejercicio}
    Si $q\in \mathbb{Q}$ es un entero algebraico, entonces $q\in \mathbb{Z}$.
\end{ejercicio}
